% =====================================================================
%  Comandos e ambientes personalizados do relatório
% =====================================================================

% ---------- Fonte abaixo de figuras e tabelas (padrão ABNT) ----------
\newcommand{\fonte}[1]{\par\vspace{3pt}{\footnotesize Fonte: #1}\par}

% ---------- Texto a ser completado pelo aluno ----------
% Aparece destacado em amarelo no PDF para não passar despercebido.
\newcommand{\preencher}[1]{\colorbox{yellow!45}{\strut\footnotesize\bfseries[PREENCHER: #1]}}

% ---------- Código/comando curto dentro do texto ----------
\newcommand{\cod}{\lstinline[basicstyle=\ttfamily\small,keywordstyle=,stringstyle=,commentstyle=,
  postbreak={},breakindent=0pt]}

% ---------- Caixa exibida quando o print ainda não foi adicionado ----------
\newcommand{\caixaprint}[2]{%
  \begin{tcolorbox}[enhanced, width=0.86\linewidth, colback=gray!4,
      frame hidden, borderline={0.7pt}{0pt}{gray!55,dashed}, arc=1.5mm,
      halign=center, valign=center, height=3.6cm, fontupper=\small,
      before skip=2pt, after skip=2pt]
    {\color{cordestaque}\bfseries INSERIR PRINT AQUI}\par\smallskip
    {\ttfamily\footnotesize Salvar como: prints/#1.png}\par\medskip
    \begin{minipage}{0.92\linewidth}\centering\footnotesize #2\end{minipage}
  \end{tcolorbox}}

% ---------- Print de tela ----------
% \printtela[largura]{arquivo}{legenda}{o que capturar}
% Se existir prints/<arquivo>.png (ou .jpg/.jpeg), a imagem é inserida;
% caso contrário, aparece a caixa indicando qual print deve ser feito.
\newcommand{\printtela}[4][0.9]{%
  \begin{figure}[H]
    \centering
    \caption{#3}\label{fig:#2}
    \IfFileExists{prints/#2.png}{\includegraphics[width=#1\linewidth]{prints/#2.png}}{%
      \IfFileExists{prints/#2.jpg}{\includegraphics[width=#1\linewidth]{prints/#2.jpg}}{%
        \IfFileExists{prints/#2.jpeg}{\includegraphics[width=#1\linewidth]{prints/#2.jpeg}}{%
          \caixaprint{#2}{#4}}}}
    \fonte{Elaborado pelo autor (2026).}
  \end{figure}}

% ---------- Código-fonte a partir de arquivo do laboratório ----------
% \codigo[opções]{linguagem}{arquivo dentro de codigos/}{legenda}{rótulo}
%   -> listagem curta: fica inteira na mesma página (dentro de uma minipage)
% \codigolongo[opções]{linguagem}{arquivo}{legenda}{rótulo}
%   -> listagem longa: pode continuar na página seguinte
\newcommand{\codigo}[5][]{%
  \par\smallskip\noindent\begin{minipage}{\linewidth}%
  \lstinputlisting[language=#2,caption={[{#4}]{#4 -- \nolinkurl{#3}}},label={cod:#5},#1]{codigos/#3}%
  \end{minipage}\par\smallskip}
\newcommand{\codigolongo}[5][]{%
  \par\needspace{12\baselineskip}%
  \lstinputlisting[language=#2,caption={[{#4}]{#4 -- \nolinkurl{#3}}},label={cod:#5},#1]{codigos/#3}}

% ---------- Listagens digitadas no texto (não quebram entre páginas) ----------
% \begin{terminal}[caption={...},label={cod:...}] ... \end{terminal}  -> comandos de terminal
% \begin{listagem}[language=...,caption={...}] ... \end{listagem}    -> outros trechos de código
\lstnewenvironment{terminal}[1][]%
  {\endgraf\smallskip\noindent\minipage{\linewidth}\lstset{style=terminal,#1}}%
  {\endminipage\endgraf\smallskip}
\lstnewenvironment{listagem}[1][]%
  {\endgraf\smallskip\noindent\minipage{\linewidth}\lstset{#1}}%
  {\endminipage\endgraf\smallskip}

% ---------- Caixas de destaque ----------
\newtcolorbox{nota}[1][Observação]{enhanced, breakable, colback=corprimaria!4,
  colframe=corprimaria, boxrule=0pt, leftrule=2.5pt, arc=0pt, outer arc=0pt,
  left=7pt, right=7pt, top=5pt, bottom=5pt, fontupper=\small,
  before upper={\textbf{\color{corprimaria}#1:}\ }, before skip=8pt, after skip=8pt}

\newtcolorbox{atencao}[1][Atenção]{enhanced, breakable, colback=cordestaque!6,
  colframe=cordestaque, boxrule=0pt, leftrule=2.5pt, arc=0pt, outer arc=0pt,
  left=7pt, right=7pt, top=5pt, bottom=5pt, fontupper=\small,
  before upper={\textbf{\color{cordestaque}#1:}\ }, before skip=8pt, after skip=8pt}

% ---------- Estilos TikZ dos diagramas ----------
\tikzset{
  caixa/.style={draw=corprimaria, fill=corprimaria!6, rounded corners=3pt, align=center,
                font=\footnotesize, inner sep=5pt, line width=0.6pt},
  caixaescura/.style={draw=corprimaria, fill=corprimaria, text=white, rounded corners=3pt,
                align=center, font=\footnotesize\bfseries, inner sep=5pt},
  caixalaranja/.style={draw=cordestaque, fill=cordestaque!8, rounded corners=3pt, align=center,
                font=\footnotesize, inner sep=5pt, line width=0.6pt},
  caixacinza/.style={draw=gray!60, fill=gray!8, rounded corners=3pt, align=center,
                font=\footnotesize, inner sep=5pt},
  seta/.style={-{Stealth[length=2.4mm]}, line width=0.7pt, color=corcinza},
  setadupla/.style={{Stealth[length=2.4mm]}-{Stealth[length=2.4mm]}, line width=0.7pt, color=corcinza},
  rotulo/.style={font=\scriptsize, text=corcinza, fill=white, inner sep=1.5pt},
}
